\begin{figure}[h]
\centering
\resizebox{\textwidth}{!}{
\begin{tikzpicture}[
    block/.style={
        rectangle,
        draw,
        rounded corners,
        align=center,
        minimum width=2.8cm,
        minimum height=1.0cm
    },
    arrow/.style={
        -{Latex},
        thick
    },
    node distance=1.4cm
]

% Main horizontal chain
\node[block] (traj) {Trajectory\\Generator};

\node[block, right=of traj] (pos)
{Position PID\\$x,y \rightarrow \phi_d,\theta_d$};

\node[block, right=of pos] (att)
{Attitude PID\\$\phi,\theta \rightarrow \Delta PWM$};

\node[block, right=of att] (mix)
{Mixer\\$PWM_1,\ldots,PWM_4$};

\node[block, right=of mix] (motor)
{Motor Model\\$PWM \rightarrow T$};

\node[block, right=of motor] (drone)
{Drone Physics\\Isaac Sim};

% Upper control blocks
\node[block, above left=1.2cm and 0.2cm of mix] (yaw)
{Yaw PID\\$\psi \rightarrow \Delta PWM_\psi$};

\node[block, above right=1.2cm and 0.2cm of mix] (alt)
{Altitude PID\\$z \rightarrow PWM_{base}$};

% Main arrows
\draw[arrow] (traj) -- (pos);
\draw[arrow] (pos) -- (att);
\draw[arrow] (att) -- (mix);
\draw[arrow] (mix) -- (motor);
\draw[arrow] (motor) -- (drone);

% Yaw and altitude arrows to mixer
\draw[arrow] (yaw) -- (mix);
\draw[arrow] (alt) -- (mix);

% Feedback line
\draw[arrow]
    (drone.south) |- ++(0,-1.0) -| (pos.south)
    node[pos=0.25, below] {state feedback};

\end{tikzpicture}
}
\caption{Cascade PID control structure}
\label{fig:cascade_pid}
\end{figure}